\hypertarget{chapter-39-advanced-template-metaprogramming}{%
\section{CHAPTER 39: ADVANCED TEMPLATE
METAPROGRAMMING}\label{chapter-39-advanced-template-metaprogramming}}

\hypertarget{advanced-template-metaprogramming}{%
\section{ADVANCED TEMPLATE
METAPROGRAMMING}\label{advanced-template-metaprogramming}}

\hypertarget{the-evolution-of-tmp}{%
\subsection{1. The Evolution of TMP}\label{the-evolution-of-tmp}}

Template Metaprogramming (TMP) is the art of using the compiler to
generate code.

\begin{itemize}
\tightlist
\item
  \textbf{C++98:} Recursive struct instantiation,
  \passthrough{\lstinline!enum!} hacks. (Hard).
\item
  \textbf{C++11:} \passthrough{\lstinline!constexpr!},
  \passthrough{\lstinline!static\_assert!},
  \passthrough{\lstinline!using!}, Variadic Templates. (Better).
\item
  \textbf{C++14:} Variable templates, \passthrough{\lstinline!auto!}
  return type. (Cleaner).
\item
  \textbf{C++17:} \passthrough{\lstinline!if constexpr!}, Fold
  expressions, \passthrough{\lstinline!void\_t!}. (Powerful).
\item
  \textbf{C++20:} Concepts (\passthrough{\lstinline!requires!}). (The
  Holy Grail).
\end{itemize}

\hypertarget{sfinae-substitution-failure-is-not-an-error}{%
\subsection{2. SFINAE (Substitution Failure Is Not An
Error)}\label{sfinae-substitution-failure-is-not-an-error}}

Before Concepts, SFINAE was the only way to constrain templates.

\hypertarget{stdenable_if}{%
\subsubsection{\texorpdfstring{2.1
\texttt{std::enable\_if}}{2.1 std::enable\_if}}\label{stdenable_if}}

\begin{lstlisting}[language={C++}]
#include <type_traits>
#include <iostream>

// Enable only for integral types
template <typename T>
typename std::enable_if<std::is_integral<T>::value, void>::type
process(T t) {
    std::cout << "Integral: " << t << "\n";
}

// Enable only for floating point
template <typename T>
typename std::enable_if<std::is_floating_point<T>::value, void>::type
process(T t) {
    std::cout << "Float: " << t << "\n";
}
\end{lstlisting}

\hypertarget{the-void_t-trick-c17}{%
\subsubsection{\texorpdfstring{2.2 The \texttt{void\_t} Trick
(C++17)}{2.2 The void\_t Trick (C++17)}}\label{the-void_t-trick-c17}}

Detecting if a type has a member function.

\begin{lstlisting}[language={C++}]
template <typename T, typename = void>
struct has_print : std::false_type {};

template <typename T>
struct has_print<T, std::void_t<decltype(std::declval<T>().print())>> : std::true_type {};

static_assert(has_print<MyClass>::value, "MyClass must have print()");
\end{lstlisting}

\hypertarget{curiously-recurring-template-pattern-crtp}{%
\subsection{3. Curiously Recurring Template Pattern
(CRTP)}\label{curiously-recurring-template-pattern-crtp}}

Static polymorphism. The base class knows the derived class type at
compile time.

\begin{lstlisting}[language={C++}]
template <typename Derived>
class Base {
public:
    void interface() {
        // Compile-time dispatch
        static_cast<Derived*>(this)->implementation();
    }
};

class Derived : public Base<Derived> {
public:
    void implementation() {
        std::cout << "Derived impl\n";
    }
};
\end{lstlisting}

\textbf{Use Case:} Mixins, adding functionality (like equality
operators) without virtual overhead.

\hypertarget{policy-based-design}{%
\subsection{4. Policy-Based Design}\label{policy-based-design}}

Designing classes that take ``policies'' (strategy classes) as template
arguments to define behavior.

\begin{lstlisting}[language={C++}]
template <typename OutputPolicy, typename LanguagePolicy>
class HelloWorld : public OutputPolicy, public LanguagePolicy {
public:
    void run() {
        print(message()); // OutputPolicy::print, LanguagePolicy::message
    }
};
\end{lstlisting}

\hypertarget{modern-tmp-with-concepts-c20}{%
\subsection{5. Modern TMP with Concepts
(C++20)}\label{modern-tmp-with-concepts-c20}}

Replacing SFINAE with readable constraints.

\begin{lstlisting}[language={C++}]
template<typename T>
concept Printable = requires(T t) {
    { t.print() } -> std::same_as<void>;
};

void process(Printable auto& obj) {
    obj.print();
}
\end{lstlisting}
